\documentclass[11pt]{article}
\usepackage{fontspec}
\usepackage[french]{babel}
\usepackage{amsmath}
\usepackage{amssymb}
\usepackage{amsthm}
\usepackage{mathrsfs}
\usepackage{tikz-cd}
\usepackage{enumitem}

\newcommand{\printedpage}[1]{\marginpar{\footnotesize p.~#1}}
\newcommand{\msp}[1]{\marginpar{\footnotesize [#1]}}
\newcommand{\Hom}{\operatorname{Hom}}
\newcommand{\Ind}{\operatorname{Ind}}
\newcommand{\Ens}{\mathrm{Ens}}
\newcommand{\Cat}{\mathrm{Cat}}
\newcommand{\Cathot}{\mathrm{Cat\,hot}}
\newcommand{\Hot}{\mathrm{Hot}}
\newcommand{\Ob}{\mathrm{Ob}}
\newcommand{\id}{\mathrm{id}}
\newcommand{\Ch}{\mathrm{Ch}}

\begin{document}

\printedpage{24}
% head: A. Grothendieck : Dérivateurs, Ch. V, Édition des Universités de Montpellier II et Paris VII
% foot: Version 9 septembre 2001

\[
\begin{tikzcd}
\Hom_C(X,K) \arrow[r,"\alpha_X"] \arrow[dr,dashed] & \Hom_{C'}(\varphi X,\varphi K) \arrow[d,dash,"\wr"',"\mathrm{adj}"] \\
 & \Hom_C(X,\psi\varphi K)
\end{tikzcd}
\]
est déduite de la flèche d'adjonction
\[
K \longrightarrow \psi\varphi K
\]
en prenant $\Hom_C(X,-)$, donc $\alpha_X$ est un isomorphisme pour tout $X$ si et seulement si cette flèche d'adjonction est un isomorphisme. Donc
\[
\text{(ii)} \Longleftrightarrow \text{(ii$'$)} \Longleftrightarrow \text{(ii$''$)}.
\]
D'autre part,
\[
\text{(ii$''$)} \Longrightarrow \text{(i)} \Longrightarrow \text{(i$'$)} \Longrightarrow \text{(i$''$)}
\]
sont tautologiques, en tenant compte pour (i) $\Longrightarrow$ (i) \texttt{[probablement (i) $\Longrightarrow$ (i$'$)]} que les flèches d'adjonction
\[
X \xrightarrow{u} \psi\varphi X \qquad\qquad\qquad \text{dans } C
\]
sont bien dans $W$, car on a un triangle commutatif
\[
\begin{tikzcd}[column sep=large]
\varphi(X) \arrow[r,"\varphi(u)"] \arrow[dr,"\id_{\varphi(X)}"] & \varphi(\psi\varphi X) = (\varphi\psi)(\varphi X) \arrow[d,"{\begin{array}{l}\wr\ \mathrm{adj}\ \varphi\psi \longrightarrow 1,\ \text{qui est un}\\ \text{isomorphisme car } \psi \text{ est pleinement fidèle}\end{array}}"] \\
 & \varphi(X),
\end{tikzcd}
\]
donc $\varphi(u)$ est bien un isomorphisme.

Reste à prouver (i$''$) $\Longrightarrow$ (ii$'$), i.e. on suppose

\noindent\msp{page 39}%
que le foncteur $X \longrightarrow \Hom(X,K)$ transforme le morphisme d'adjonction particulier
\[
K \xrightarrow{i} \psi\varphi K
\]
en
\[
\begin{aligned}
u \quad &\longrightarrow \quad u\circ i\\
\Hom(\psi\varphi(K),K) \quad &\longrightarrow \quad \Hom(K,K),
\end{aligned}
\]
qui est une bijection. Donc $\id_K$ est dans l'image, soit $p$ un élément du premier membre qui lui donne naissance

\printedpage{25}
% head: A. Grothendieck : Dérivateurs, Ch. V, Édition des Universités de Montpellier II et Paris VII
% foot: Version 9 septembre 2001

\[
\begin{tikzcd}
K \arrow[r,hook,shift left,"i"] & \psi\,\underbrace{\varphi(K)}_{K'} \arrow[l,shift left,"p"]
\end{tikzcd}
\qquad\qquad {\scriptstyle \text{donc } p\circ i = \id_K}
\]
Considérons
\[
ip = q : \psi(K') \longrightarrow \psi(K') \qquad\qquad\qquad (\text{NB } q^2 = q),
\]
comme $\psi$ est pleinement fidèle, il existe un unique
\[
q_0 : K' \longrightarrow K' \qquad\qquad\qquad (K' \overset{\text{déf}}{=} \varphi(K))
\]
tel que l'on ait
\[
q = \psi(q_0).
\]
La relation $pi = \id_K$ donne, en appliquant $\varphi$,
\[
\varphi(p)\varphi(i) = \id,
\]
donc comme $\varphi(i)$ est un isomorphisme (c'est
\[
\begin{tikzcd}[column sep=large]
\varphi(K) \arrow[r,"\sim"] \arrow[dr,"\sim"{sloped,pos=0.45},"\id"{pos=0.6}] & \varphi(\psi\varphi)(K) = (\varphi\psi)(\varphi K) \arrow[d,"\wr"',"\mathrm{adj}"] \\
 & \varphi(K),
\end{tikzcd}
\]
rendant commutatif le triangle ci-contre), $\varphi(p)$ l'est. Donc $\varphi(q) = \varphi(i)\varphi(p)$ l'est, puisque $\varphi(i)$ et $\varphi(i)$ \texttt{[plutôt $\varphi(p)$]} le sont. Mais
\[
\varphi(q) = \varphi\psi(q_0) \quad \underbrace{\xrightarrow{\ \sim\ }}_{\mathrm{adj}\ \varphi\psi\ \xrightarrow{\sim}\ \id} \quad q_0,
\]
donc $q_0$ est un isomorphisme, donc aussi $q = \psi q_0$. Comme $q^2 = q$, cela implique $q = \id$. Donc $i$ et $p$ sont des isomorphismes inverses l'un de l'autre, ce qui établit (ii$''$) \texttt{[plutôt (ii$'$)]} et achève la démonstration.

\textbf{Corollaire :} \textit{Considérons des foncteurs $C^o \xrightarrow{H} \Ens$, et $C'^o \xrightarrow{H'} \Ens$, tels que $H \simeq H'\circ\varphi$. Pour que $H$ soit représentable, il faut et il suffit que $H'$ le soit, et dans ce cas, l'objet $K$ de $C$}

\noindent\msp{page 40}%
\textit{qui représente $H$ satisfait aux conditions équivalentes de la proposition. Si $K$, $K'$ les représentent, on a des isomorphismes canoniques $K' \simeq \varphi K$, $K \simeq \psi K'$.}

Si $H'$ est représentable par $K' \in \Ob C'$, on aura donc
\[
H'\circ\varphi(X) \simeq \Hom_{C'}(\varphi X,K') \overset{\mathrm{adj}}{\simeq} \Hom_C(X,\psi(K'))
\]
\printedpage{26}
% head: A. Grothendieck : Dérivateurs, Ch. V, Édition des Universités de Montpellier II et Paris VII
% foot: Version 9 septembre 2001
donc $H$ est représentable par $K = \psi(K')$, satisfaisant la condition (ii) de la proposition. Inversement, si $H$ est représentable par un objet $K$, alors $K$ satisfait à la condition (i), donc par (ii$''$) on a
\[
\Hom_C(X,K) \xrightarrow{\sim} \Hom_{C'}(\varphi X,\varphi K),
\]
donc le foncteur $H' : C^o \longrightarrow \Ens$ qui factorise $H$ s'explicite comme $X' \longmapsto \Hom_{C'}(X',\varphi K)$, donc est représentable par $\varphi K$. % [sic] C^o -> C'^o, H' is a functor on C'^o (cf. the Corollaire)

Passons aux ind-objets. Les foncteurs $\varphi$, $\psi$ passent aux ind-objets, d'où des foncteurs
\[
\begin{tikzcd}[column sep=large]
\Ind C \arrow[r,shift left,"\Psi = \Ind\psi"] \arrow[r,shift right,"\Phi = \Ind\varphi"'] & \Ind C'.
\end{tikzcd} % [sic] both arrows printed pointing to Ind C' -> the arrow Psi = Ind psi apparently intended from Ind C' to Ind C, psi goes from C' to C
\]
On aura encore la formule d'adjonction pour $X = (X_i)_{i\in I}$ dans $\Ind C$, $X' = (X'_{i'})_{i'\in I'}$ dans $\Ind C'$, car
\[
\begin{aligned}
\Hom(X,\Psi X') &= \varprojlim_i \varinjlim_{i'} \Hom(X_i,\psi(X'_{i'}))\\
&= \varprojlim_i \varinjlim_{i'} \Hom(\varphi X_i,X'_{i'})\\
&= \Hom(\Phi X,X').
\end{aligned}
\]

\noindent\msp{page 41}%
On prouve de même que $\Psi$ est pleinement fidèle, car si $X' = (X'_i)_{i\in I}$, $Y' = (Y'_j)_{j\in J}$ sont des objets de $\Ind C'$, on aura
\[
\begin{tikzcd}[column sep=small]
\Hom(X',Y') \arrow[d,"\Psi_{X',Y'}"] & \simeq \varprojlim\limits_i \varinjlim\limits_j & \Hom_{C'}(X'_i,Y'_j) \arrow[d,"\psi_{X'_i,Y'_j}"] \\
\Hom(\Psi X',\Psi Y') & \simeq \varprojlim\limits_i \varinjlim\limits_j & \Hom_C(\psi(X'_i),\psi(Y'_j)),
\end{tikzcd}
\]
donc $\Psi_{X',Y'}$ est un isomorphisme si les $\psi_{X'_i,Y'_j}$ le sont.

Donc la proposition 1 peut s'appliquer à la situation $\Phi$, $\Psi$ entre ind-objets. Je vais l'expliciter sous la forme la plus commode :

\textbf{Proposition 2.} \textit{Soit $K = (K_i)_{i\in I}$ un ind-objet de $C$, et soit $H$ le foncteur}
\[
\left\{
\begin{aligned}
H : X \quad &\longmapsto \quad \underbrace{\Hom_C(X,K)}_{K(X)} \overset{\text{déf}}{=} \varinjlim_i \Hom_C(X,K_i)\\
C^o \quad &\longrightarrow \quad \Ens
\end{aligned}
\right.
\]
\textit{qu'il ind-représente. Conditions équivalentes}

\begin{enumerate}[label={}]
\item[(i)] \textit{Le foncteur $H$ se factorise (à isomorphisme près) par $C'$ via $\varphi$, i.e. transforme les $u \in W_\varphi$ en isomorphismes.}
\printedpage{27}
% head: A. Grothendieck : Dérivateurs, Ch. V, Édition des Universités de Montpellier II et Paris VII
% foot: Version 9 septembre 2001
\item[(i$'$)] \textit{Même condition pour les flèches d'adjonction $X \longrightarrow \psi\varphi(X)$ (lesquelles sont dans $W$).}
\item[(i$''$)] \textit{Même condition pour la flèche d'adjonction $K \longrightarrow \psi\varphi(K)$.}
\item[(ii)] \textit{$K$ appartient à l'image essentielle de} $\Ind C'$ \textit{par $\Psi$.}
\item[(ii$'$)] $K \longrightarrow \Psi\Phi(K)$ \textit{est un} isomorphisme \textit{dans} $\Ind C$.) % [sic] ".)" -> ".", closing parenthesis without opening parenthesis
\item[(ii$''$)] \textit{L'homomorphisme fonctoriel}
\[
K(X) = \Hom_C(X,K) \longrightarrow \Hom_{C'}(\varphi(X),\varphi(K))
\]
\texttt{[plutôt} $\Hom_{\mathrm{Ind}(C)}(X,K) \longrightarrow \Hom_{\mathrm{Ind}(C')}(\Phi(X),\Phi(K))$\texttt{, en identifiant l'objet} $X$ \texttt{de} $C$ \texttt{à son image dans} $\mathrm{Ind}(C)$\texttt{]} \textit{est un isomorphisme.}
\end{enumerate}

\textsc{Démonstration.} L'équivalence de (ii), (ii$'$), (ii$''$) est la proposition 1 appliquée à $\Phi$ et $\Psi$, à cela près qu'il faut noter que (ii$''$) reste évidemment valable si on y remplace $X$ par un ind-objet $(X_j)_{j\in J}$ quelconque. On a

\noindent\msp{page 42}%
\[
\text{(ii$''$)} \Longrightarrow \text{(i)} \Longrightarrow \text{(i$'$)}
\]
tautologiquement, il reste à prouver que (i$'$) implique (i$''$) et (i$''$) (ii$'$). On ne peut pas appliquer directement la proposition 1, car l'hypothèse (i$'$) ici ne porte que sur les objets de $\Ind C$ qui proviennent de $C$, et il n'est pas évident à priori qu'elle passe aux objets quelconques. Mais cela l'est dans le cas d'un homomorphisme de ind-objets $(X_i)$, $(Y_i)$ sur le même ensemble d'indices, donné par un homomorphisme ordinaire (terme à terme) de systèmes inductifs sur $I$, $X_i \longrightarrow Y_i$: si les $X_i \longrightarrow Y_i$ sont dans $W$, l'hypothèse (i) implique que $H(Y) \longrightarrow H(X)$ est un isomorphisme, et de même pour (i$'$) pour $\underbrace{X}_{=(X_i)_i} \;\longrightarrow\; \underbrace{\Psi\Phi(X)}_{=(\psi\varphi(X_i))_i}$\,, où l'homomorphisme d'adjonction est donné terme à terme par les homomorphismes d'adjonction $X_i \longrightarrow \psi\varphi(X_i)$. Donc l'hypothèse (i$'$) ici implique l'hypothèse (i$'$) de la proposition 1 pour $\Phi$ et $\Psi$, donc aussi (ii$'$) (par la proposition 1), et on gagne. % [sic] "et (i'') (ii')" -> "et (i'') => (ii')", the implication sign or verb is missing

Revenons à la situation où $C = \Cathot$, $C' = \Hot$, $\varphi$ le foncteur \texttt{[de]} localisation. Si $K \in \Ob\Cathot$, alors $\underline{\psi\varphi(K)}$ s'appelle \textit{l'enveloppe de Thomason-Kan} de $K$ - si $K$ est une catégorie, son enveloppe $\bar{K}$ et la flèche $K \longrightarrow \bar{K}$ ne sont définis que dans la catégorie Cat\,hot, i.e. $\bar{K}$ n'est défini qu'à `homotopisme près'. % ?? uncertain: underline under psi varphi(K) -- underlined / not underlined

\noindent\msp{page 43}%
Mais on peut toujours \textit{choisir} un objet $\bar{K}$ de Cat, et une flèche
\[
i : K \longrightarrow \bar{K} \qquad\qquad\qquad \textit{dans } \Cat,
\]
\printedpage{28}
% head: A. Grothendieck : Dérivateurs, Ch. V, Édition des Universités de Montpellier II et Paris VII
% foot: Version 9 septembre 2001
pas seulement dans Cat\,hot, de telle façon de plus que $i$ soit une `\textit{cofibration} de Thomason' (et en particulier une immersion ouverte?). Il en résulte que si on a un système inductif
\[
K = (K_n)_{n\in\mathbf{N}}
\]
indexé par $\mathbf{N}$, et si on choisit pour chaque $n$ une enveloppe de Thomason-Kan
\[
K_n \overset{\alpha_n}{\hookrightarrow} \bar{K}_n
\]
de cette façon, on peut trouver des foncteurs de transition
\[
\bar{K}_n \xrightarrow{\bar{i}_n} \bar{K}_{n+1}
\]
tels que les carrés
\[
\begin{tikzcd}
K_n \arrow[r,hook,"\alpha_n"] \arrow[d,"i_n"'] & \bar{K}_n \arrow[d,"\bar{i}_n"] \\
K_{n+1} \arrow[r,hook,"\alpha_{n+1}"] & \bar{K}_{n+1}
\end{tikzcd}
\]
soient commutatifs, de telle sorte que l'on ait donc un homomorphisme de \textit{systèmes inductifs de} Cat
\[
K \longrightarrow \bar{K}
\]
et a fortiori un homomorphisme de ind-objets de Cat. Ceci posé, les conditions (i), (ii$'$) de la proposition 2 montrent que le foncteur ind-représenté par $K$ \textit{dans} $C = \Cathot$ transforme quasi-isomorphismes

\noindent\msp{page 44}%
ordinaires (équivalences faibles) en bijections, \textit{si et seulement si} l'homomorphisme précédent est un isomorphisme de ind-objets de Ind\,Cat\,hot.

Ce sera donc cela la condition naturelle à imposer aux ind-catégories de chemins
\[
\underline{\underline{\Ch}}_{\Sigma}(\mathcal{M};a,b) \qquad\qquad\qquad (a,b \in \mathrm{ob}\,\mathcal{M})
\]
provenant d'une `catégorie de modèles' $\mathcal{M}$, $\Sigma$.

Bien sûr, disposant d'un système inductif explicite de catégories (pas seulement d'objets de Cat\,hot)
\[
K_n = \underline{\Ch}_{\Sigma,n}(\mathcal{M};a,b),
\]
il est tentant de passer à la limite inductive
\[
K_\infty = \varinjlim_n \underline{\Ch}_{\Sigma,n}(\mathcal{M};a,b) \overset{\text{déf}}{=} \underline{\Ch}_{\Sigma}(\mathcal{M};a,b),
\]
et ceci d'autant plus que les foncteurs de transition sont des monomorphismes et pleinement fidèles, de sorte que les $K_n$ s'identifient à une suite de sous-catégories pleines de

\end{document}
